\documentclass[conference]{IEEEtran}
\IEEEoverridecommandlockouts
% The preceding line is only needed to identify funding in the first footnote. If that is unneeded, please comment it out.
\usepackage{cite}
\usepackage{amsmath,amssymb,amsfonts}
\usepackage{algorithmic}
\usepackage{graphicx}
\usepackage{textcomp}
\usepackage{xcolor}
\def\BibTeX{{\rm B\kern-.05em{\sc i\kern-.025em b}\kern-.08em
    T\kern-.1667em\lower.7ex\hbox{E}\kern-.125emX}}
\begin{document}

\title{Moiré Potential Landscapes Dictate the Thermal Cycling Stability of Twist-Angle-Engineered Thermal Anisotropy in Van Der Waals Heterostructures}\author{Anonymous Research Plan Author}\maketitle

\begin{abstract}
Twist-angle engineering in van der Waals heterostructures establishes a powerful degree of freedom for thermal management. In random-stacked transition metal dichalcogenide films, interlayer rotation impedes through-plane transport while preserving in-plane crystallinity, producing a room-temperature thermal anisotropy ratio of approximately 900. This extreme decoupling arises from rotational disorder that breaks cross-plane translational symmetry, inducing glass-like phonon transport without expanding the interlayer spacing. However, the structural stability of these engineered configurations under sustained thermal cycling remains unaddressed. The existing evidence validates the anisotropy mechanism only for MoS2 and WS2, leaving both the material generality and the operational durability of twist-angle-engineered thermal anisotropy unresolved.

This proposal replaces the assumed shear-strain locking mechanism with a moiré-potential landscape model, positing that twist-angle-dependent energy barriers dictate the critical interlayer binding threshold for structural stability. A factorial design-of-experiments will measure cross-plane thermal conductivity and twist-angle retention before and after controlled thermal cycling across diverse van der Waals material families. The design discriminates the moiré mechanism from defect-mediated or chemical pinning by comparing retention across materials with varying moiré potential depths but similar intrinsic defect densities. A null or contradictory outcome branch triggers an audit of construct validity and a revision of the boundary claim, while an uninformative branch halts interpretation until measurement prerequisites are confirmed.
\end{abstract}
\begin{IEEEkeywords}
van der Waals heterostructures, interlayer twist angle, moiré potential landscape, thermal cycling stability, cross-plane thermal conductivity, interlayer binding energy
\end{IEEEkeywords}\section{Introduction}
\subsection{Thermal Anisotropy and the Structural Stability Gap}
Van der Waals heterostructures offer a scalable platform for manipulating heat flow by decoupling in-plane and cross-plane thermal transport. Random interlayer rotations in transition metal dichalcogenide films impede through-plane phonon propagation while preserving in-plane crystallinity, yielding a room-temperature thermal anisotropy ratio of approximately 900 (W3202213959). This extreme anisotropy arises because rotational disorder breaks the translational symmetry required for efficient transverse phonon coupling, effectively rendering the cross-plane direction thermally insulating. However, the evidence establishing interlayer rotation as a design degree of freedom is bounded: the mechanism is validated only for MoS2 and WS2, and the persistence of these rotational configurations under sustained thermal cycling and operational stress gradients remains unaddressed. A proposal that treats twist-angle-engineered thermal anisotropy as a durable device feature must therefore account for whether the engineered stacking registry survives the very thermal gradients it is designed to mitigate.~\cite{cite_26ecaf73d88dfe33}

\subsection{Mechanism Replacement: Moiré Potential as the Dominant Pinning Source}
The selected research direction replaces the assumed shear-strain locking mechanism with a twist-angle-dependent moiré potential landscape as the fundamental source of structural stability. The proposed contribution posits that engineered interlayer twist angles remain stable under thermal cycling provided the interlayer binding energy exceeds a critical threshold determined by local moiré potential energy barriers, and provided that defect-mediated or chemically-induced interlayer pinning does not dominate. This mechanism-replacement framing is necessary because alternative explanations are not negligible: thermal cycling can induce interlayer diffusion or chemical bonding that pins layers independently of the moiré potential, and defect-mediated pinning via chalcogen vacancies could maintain suppressed cross-plane thermal conductivity while masking the true moiré threshold. Distinguishing the moiré mechanism from these confounders requires a factorial design that varies material family, twist angle, and cycling intervals while controlling for defect density and interfacial contact resistance changes.

\subsection{Scope, Design Consequences, and Transition}
The claim scope is explicitly bounded to regimes free from defect-mediated or chemically-induced pinning. This boundary is not a caveat but a constitutive condition of the proposed mechanism: if interlayer binding energy falls below the moiré-determined critical threshold, or if chemical bonding becomes the dominant retention force, the stability conclusion fails. The design consequence is a pre-cycling baseline measurement followed by repeated structural and thermal characterization after defined cycling intervals, with twist-angle retention assessed against the pre-cycling uncertainty interval. The frozen upstream evidence supports the measurability of the relevant observables---through-plane conductivity via time-domain thermoreflectance and layer-resolved structural modes via Raman spectroscopy---but the numerical targets for the twist-angle grid, cycling protocol, and decision thresholds require human specification. This section establishes the mechanism-replacement hypothesis and its boundary conditions; the next stage operationalizes the factorial design that can discriminate the moiré potential mechanism from its confounders.~\cite{cite_26ecaf73d88dfe33,cite_c8ce7e25ee0b4157}

\section{Background, Survey, and Research Gap}
\subsection{The Interface-Dominated Thermal Transport Regime}
The relentless miniaturization of electronic and energy conversion devices has elevated thermal management from a secondary engineering consideration to a primary bottleneck governing reliability and speed. As feature sizes shrink into the nanoscale regime, spatially concentrated power generation exposes the breakdown of classical continuum assumptions, necessitating a rigorous examination of the fundamental physical limits inherent to charge and heat carrier dynamics. A central trend emerging across contemporary research is the paradigm shift toward interface-dominated regimes, where thermal boundary conductance supersedes bulk conduction, effectively dictating the ultimate heat dissipation capabilities of integrated systems. This transition demands a move beyond passive material optimization toward active interface engineering, where the chemical and physical origins of boundary conductance are manipulated to control heat flow.~\cite{cite_c8ce7e25ee0b4157}

\subsection{Anisotropy and Interlayer Rotation as Design Degrees of Freedom}
In van der Waals (vdW) heterostructures, the stark dichotomy between strong in-plane covalent bonding and weak out-of-plane vdW interactions naturally leads to highly anisotropic thermal conductivities. Recent advances demonstrate that structural engineering can transcend natural limits by decoupling in-plane and cross-plane transport mechanisms. Specifically, the realization of extreme anisotropy through interlayer rotational disorder has emerged as a pivotal development. By stacking monolayers with random crystalline orientations, researchers have constructed polycrystalline transition metal dichalcogenide (TMD) films that exhibit ultralow through-plane thermal conductivity while preserving high in-plane thermal transport. This approach yields thermal anisotropy ratios reaching approximately 900 for MoS2, surpassing natural anisotropic benchmarks such as pyrolytic graphite. The through-plane conductivity drops to values comparable to that of air, effectively rendering the cross-plane direction thermally insulating, while the in-plane conductivity remains robust.~\cite{cite_26ecaf73d88dfe33}

\subsection{The Material Generality Boundary}
Despite the demonstrated efficacy of random interlayer rotations in suppressing cross-plane thermal transport, the underlying mechanism remains narrowly validated. Current evidence confirms this extreme thermal anisotropy mechanism exclusively for MoS2 and WS2. This restriction leaves the applicability to other van der Waals materials unverified. The survey highlights that it remains unknown whether random interlayer rotations yield similar anisotropy ratios in materials with different phonon dispersion characteristics or interlayer bonding strengths. Without addressing this material generality, the design principle of using twist angles to engineer thermal anisotropy cannot be generalized across the broader vdW material family, limiting its utility for diverse thermal management applications.~\cite{cite_26ecaf73d88dfe33}

\subsection{The Structural Stability Under Operation Gap}
A more critical unresolved bridge concerns the persistence of these engineered rotational configurations under operational conditions. While interlayer rotation is established as a design degree of freedom for directed heat transport, the stability of these configurations under sustained thermal cycling and operational stress gradients is unaddressed. The current literature does not discuss whether these rotational states remain stable or relax over time when subjected to the very thermal gradients they are engineered to mitigate. This missing assumption regarding structural stability under operation prevents the translation of laboratory-scale anisotropy into reliable, long-term device performance.~\cite{cite_26ecaf73d88dfe33}

\subsection{Competing Mechanisms and Confounding Variables}
The gap in structural stability is compounded by the presence of competing retention mechanisms that have not been disentangled from the intrinsic moiré potential landscape. Thermal cycling can induce interlayer diffusion or chemical bonding that pins the layers independently of the moiré potential. Similarly, defect-mediated interlayer pinning, such as via chalcogen vacancies, or thermally induced chemical bonding could maintain suppressed cross-plane thermal conductivity, mimicking the effect of moiré potential energy barriers. Furthermore, strain relaxation via wrinkle formation or localized delamination could alter thermal transport without actual twist-angle relaxation. The inability to distinguish the true moiré potential threshold from these confounding defect-mediated or chemically-induced pinning mechanisms constitutes a significant barrier to establishing a predictive design framework for twist-angle-engineered thermal anisotropy.~\cite{cite_b9cc1fcb55e94349}

\subsection{Transition to Formal Argument}
Consequently, the research plan must pivot from merely observing anisotropy to establishing the evidence-bounded gap that motivates a formal mechanistic argument. The proposed contribution addresses this by positing that the twist-angle-dependent moiré potential landscape creates fundamental energy barriers that pin specific twist angles, dictating the critical threshold for interlayer binding energy and structural stability under thermal stress. This mechanism is asserted to be distinct from and dominant over defect-mediated or chemical pinning mechanisms, provided specific boundary conditions are met. The following sections will formalize this hypothesis, defining the critical thresholds and designing the experimental protocols necessary to isolate the moiré potential mechanism from confounding variables.~\cite{cite_6b49ace365c37cc5}

\section{Research Questions and Planned Contributions}
\subsection{Background and Operational Research Questions}
The survey establishes that engineered interlayer rotation in transition metal dichalcogenide films suppresses cross-plane thermal transport while preserving in-plane crystallinity, yielding a room-temperature thermal anisotropy ratio of approximately 900. This mechanism, however, is validated exclusively for MoS2 and WS2, and the literature does not address whether these rotational configurations remain stable under sustained thermal cycling and operational stress gradients. The proposed research replaces the assumed shear-strain locking mechanism with a twist-angle-dependent moiré potential landscape, positing that fundamental energy barriers pin specific twist angles and dictate the critical threshold for structural stability. This route operationalizes three research questions. First, which van der Waals material families exhibit moiré potential depths sufficient to maintain suppressed cross-plane conductivity after repeated thermal excursions? Second, what numerical targets for the interlayer twist-angle grid and thermal-cycling protocol isolate the moiré mechanism from defect-mediated or chemically-induced interlayer pinning? Third, how does the retention of cross-plane thermal anisotropy scale with interlayer binding energy across diverse material systems? These questions define the scope of the design, transitioning from the observation of static anisotropy to the evaluation of dynamic structural stability under operational thermal stress.~\cite{cite_26ecaf73d88dfe33,cite_6b49ace365c37cc5}

\subsection{Planned Contributions and Distinct Roles}
\begin{itemize}
\item The planned contributions map the unresolved design parameters to specific experimental and analytical roles. First, the study will construct a factorial design-of-experiments library of van der Waals heterostructures, comparing material families with varying moiré potential depths but similar intrinsic defect densities. This ablation isolates the moiré potential mechanism from confounding interlayer diffusion or chemical bonding. Second, the design will establish a measurement protocol correlating pre- and post-cycling cross-plane thermal conductivity with structural integrity scores. This contribution defines the operational boundary conditions for thermal anisotropy retention. Third, the study will develop a decision rule for twist-angle stability, evaluating whether measured angles remain within the pre-cycling uncertainty interval. Each contribution targets a distinct gap: material generality, structural stability under operation, and the isolation of the dominant retention mechanism.
\end{itemize}~\cite{cite_26ecaf73d88dfe33}

\begin{table*}[htbp]
\centering
\footnotesize
\setlength{\tabcolsep}{3pt}
\renewcommand{\arraystretch}{1.12}
\caption{Unresolved Items Delimiting the Questions}
\label{tab:research-questions-and-contributions-block-003}
\begin{tabular}{|p{0.23\textwidth}|p{0.23\textwidth}|p{0.23\textwidth}|p{0.23\textwidth}|}
\hline
Design Parameter & Operational Role & Status & Required Input \\\hline
Material Families & Determines the generality of the moiré potential mechanism. & needs\_human\_input & Specify van der Waals material families for comparison. \\\hline
Twist-Angle Grid & Defines the parameter space for rotational stability. & needs\_human\_input & Provide numerical targets for angle grid and cycling protocol. \\\hline
Sample Eligibility & Ensures measurable cross-plane thermal conductivity. & needs\_human\_input & Define eligibility criteria and minimum measurable areas. \\\hline
Calibration Standards & Validates structural and thermal measurements. & needs\_human\_input & Specify instrument calibration and quality-control procedures. \\\hline
Replication Structure & Establishes statistical power for stability claims. & needs\_human\_input & Define replication and batch structure for synthesis and cycling. \\\hline
Statistical Model & Determines the threshold for hypothesis support. & needs\_human\_input & Specify statistical model and decision thresholds. \\\hline
\end{tabular}
\end{table*}~\cite{cite_26ecaf73d88dfe33}

\section{Idea Source Checkpoints and Direction Selection Audit}
\subsection{Scientific Attraction and Evidence Anchoring}
The proposal originates from a precise evidentiary deficit in the literature on van der Waals (vdW) thermal conductors. Existing studies demonstrate that random interlayer rotations in transition metal dichalcogenide (TMD) films impede through-plane thermal transport while preserving in-plane crystallinity, yielding room-temperature thermal anisotropy ratios of approximately 900. However, this mechanism is validated exclusively for MoS2 and WS2, leaving its applicability to other vdW material families unverified. Crucially, the literature does not address the persistence of these rotational configurations under sustained thermal cycling and operational stress gradients---the very conditions these materials are engineered to mitigate. The selected direction, "Moiré Potential Landscapes Dictate the Thermal Cycling Stability," advances the proposal by replacing the assumed shear-strain locking mechanism with a twist-angle-dependent moiré potential landscape. This shift posits that fundamental energy barriers pin specific twist angles, dictating a critical threshold for interlayer binding energy that governs structural stability under thermal stress.~\cite{cite_26ecaf73d88dfe33}

\subsection{Defects Exposed by Checkpoints and Mechanism Replacement}
The audit snapshots expose two critical defects in the initial hypothesis: the unverified assumption that interlayer binding energies can be accurately calculated across diverse material families, and the confounding influence of defect-mediated or chemically-induced interlayer pinning. Thermal cycling can induce interlayer diffusion or chalcogen vacancies that pin layers independently of the moiré potential, masking the true energy threshold. To address these unsupported links, the retained direction explicitly bounds the claim scope to regimes free from defect-mediated pinning. The mechanism is redefined: the moiré potential landscape creates energy barriers that dominate retention, distinct from and superior to defect or chemical pinning. This requires a discriminating observation comparing retention across material families with varying moiré potential depths but similar intrinsic defect densities, or utilizing low-temperature cycling to suppress chemical bonding. The proposal thus shifts from a general stability claim to a conditional mechanism-replacement argument, isolating the moiré potential as the primary determinant of thermal anisotropy persistence.

\subsection{Qualified Retained Direction and Explicit Exclusions}
The qualified retained direction asserts that engineered twist angles remain structurally stable under thermal cycling provided the interlayer binding energy exceeds a critical threshold determined by the local moiré potential, and in the absence of confounding pinning mechanisms. This claim is explicitly bounded by the exclusion of regimes where defect-mediated pinning or thermally induced chemical bonding become dominant. The structural stability fails when binding energy falls below the critical threshold, or when alternative retention mechanisms mask the moiré potential's effect. This boundary condition is not a limitation but a defining feature of the proposed mechanism, distinguishing it from prior work that assumed rotational disorder was static. The transition to the next stage requires the definition of operational thermal cycling parameters and the specification of material families that allow for the isolation of moiré potential effects from defect chemistry, ensuring the experimental design can adjudicate between the proposed mechanism and the excluded alternatives.

\section{Problem Definition, Assumptions, and Hypotheses}
\subsection{Problem Statement and Admissible Domain}
Twist-angle engineering in van der Waals heterostructures establishes a design degree of freedom for directed heat transport, but the literature leaves the persistence of these rotational configurations under sustained thermal cycling and operational stress gradients unaddressed. The extreme anisotropy mechanism is validated only for MoS2 and WS2, leaving applicability to other van der Waals materials unverified. This proposal defines the formal problem: determining whether engineered interlayer twist angles remain structurally stable, and whether the twist-angle-dependent moiré potential landscape creates fundamental energy barriers that pin specific angles, distinct from and dominant over defect-mediated or chemical pinning. The admissible domain comprises heterostructure stacks of two or more 2D layers transferred to a designated substrate with a target twist angle set during assembly, subjected to repeated controlled temperature excursions within a declared range and number of cycles. The domain excludes regimes where defect-mediated pinning or interlayer chemical bonding becomes the dominant retention mechanism.~\cite{cite_26ecaf73d88dfe33}

\subsection{Core Definitions and Operational Variables}
\paragraph{Definition.} The formal object is the twist-angle-dependent moiré potential landscape, which defines the energy barriers governing interlayer rotational stability. Three operational variables anchor the problem. First, twist-angle retention is the difference between the nominal assembly angle and the angle measured on the same region of the sample after a defined cycling interval. Second, thermal cycling is a repeated sequence of controlled temperature excursions and returns within a declared range and number of cycles. Third, cross-plane thermal conductivity is measured via time-domain thermoreflectance, which requires a laser system with a 785 nm wavelength, 74.86 MHz repetition rate, and 18 mW total power, while maintaining a steady-state temperature rise of less than 4 K. In random stacked TMD films, interlayer rotation impedes through-plane thermal transport while maintaining in-plane crystallinity, leading to a room-temperature thermal anisotropy ratio of approximately 900. These definitions fix the measurable quantities that later derivations and decision protocols consume.~\cite{cite_26ecaf73d88dfe33}

\subsection{Candidate Hypothesis H1 and Mechanism}
\paragraph{Proposition (Candidate).} Hypothesis H1 (proposed, unverified): Engineered interlayer twist angles designed to suppress cross-plane thermal conductivity in van der Waals heterostructures remain structurally stable under sustained thermal cycling and operational stress gradients, provided the interlayer binding energy exceeds a critical threshold determined by the twist-angle-dependent moiré potential landscape, and in the absence of confounding defect-mediated or chemically-induced interlayer pinning. The mechanism asserts that the moiré potential landscape creates fundamental energy barriers that pin specific twist angles, dictating the critical threshold for interlayer binding energy and structural stability under thermal stress, distinct from and dominant over defect-mediated or chemical pinning mechanisms. When the interlayer binding energy exceeds the thermal activation barrier at operational temperatures, and defect-mediated or chemical pinning is negligible, the engineered twist angles persist, maintaining suppressed cross-plane thermal conductivity; otherwise, rotational relaxation degrades the thermal bottleneck.

\subsection{Numbered Problem Relation}
kappa\_\{\textbackslash{}perp\}(\textbackslash{}theta, N\_\{c\}) = f(E\_\{b\}(\textbackslash{}theta), k\_\{B\}T, \textbackslash{}Delta\textbackslash{}theta(N\_\{c\})) \textbackslash{}cdot g(\textbackslash{}kappa\_\{\textbackslash{}parallel\}, \textbackslash{}kappa\_\{int\})

\subsection{Equation Interpretation and Assumption Ledger}
The relation states that cross-plane thermal conductivity kappa\_\{\textbackslash{}perp\} depends on the interlayer binding energy E\_\{b\} at twist angle \textbackslash{}theta, the thermal energy k\_\{B\}T, and the twist-angle drift \textbackslash{}Delta\textbackslash{}theta accumulated over N\_\{c\} cycles, modulated by in-plane conductivity \textbackslash{}kappa\_\{\textbackslash{}parallel\} and interfacial conductance \textbackslash{}kappa\_\{int\}. This is a proposed functional form, not a derived law; the exact equation or boundary condition for the critical threshold is not provided and requires human input. The assumption ledger includes: (1) interlayer binding energy and moiré potential depths can be accurately calculated or measured for diverse van der Waals material families; (2) thermal cycling conditions can be selected to minimize interlayer diffusion and chemical bonding, ensuring moiré potential pinning is the dominant retention mechanism. Both assumptions are design-level premises, not verified facts. The specific synthesis and measurement protocol for controlled interlayer twist angles, the measurement technique for intrinsic cross-plane thermal conductivity, and the parameters for sustained thermal cycling (frequency, amplitude, duration) are unresolved and flagged for human confirmation.

\subsection{Failure Condition and Transition}
The declared failure condition states that structural stability of the engineered twist angles fails when the interlayer binding energy falls below the critical threshold determined by the local moiré potential energy barriers, or when defect-mediated pinning and interlayer chemical bonding become the dominant retention mechanisms, masking the true moiré potential threshold. A would-falsify condition is the demonstration that engineered twist angles relax to equilibrium stacking within device-relevant thermal cycling conditions in the absence of defect-mediated or chemical pinning, resulting in recovery of cross-plane conductivity to within 20\% of pristine aligned values. The decision rule accepts structural stability if the twist angle, measured after thermal cycling, remains within the pre-cycling uncertainty interval and the cross-plane thermal conductivity retention exceeds the pre-specified threshold. These scoped premises, the admissible domain, and the assumption ledger supply the formal problem that the expected outcomes stage consumes to map prespecified outcome branches to allowed conclusions and next actions.

\section{Study Design and Methods}
\subsection{Unit of Analysis and Source-Bounded Protocol}
\paragraph{Planned protocol.} The unit of analysis is a single van der Waals heterostructure stack with a defined interlayer twist angle. The protocol proceeds in three phases: (1) baseline characterization of the nominal twist angle and cross-plane thermal conductivity prior to cycling; (2) exposure to a declared sequence of controlled temperature excursions and returns; and (3) repeated measurement of twist-angle retention and thermal anisotropy after defined cycling intervals. The experimental unit is the stack itself, not the ensemble of transferred flakes; each stack is measured at the same spatial region across all intervals to isolate intrinsic relaxation from inter-sample variability.~\cite{cite_26ecaf73d88dfe33,cite_7a8cfa26940e2a8f}

\begin{table*}[htbp]
\centering
\footnotesize
\setlength{\tabcolsep}{3pt}
\renewcommand{\arraystretch}{1.12}
\caption{Comparators, Ablations, and Robustness Checks}
\label{tab:study-design-and-methods-b2-comparison}
\begin{tabular}{|p{0.184\textwidth}|p{0.184\textwidth}|p{0.184\textwidth}|p{0.184\textwidth}|p{0.184\textwidth}|}
\hline
Group & Twist Angle & Material Family & Cycling Regime & Purpose \\\hline
Primary & Declared grid & MoS2/WS2 & Sustained & Establish baseline retention \\\hline
Material Ablation & Declared grid & BN, black phosphorus, metal thio(seleno)phosphates & Sustained & Test material generality beyond MoS2/WS2 \\\hline
Chemical-Pinning Control & Declared grid & MoS2/WS2 & Low-temperature cycling & Suppress interlayer diffusion and chemical bonding to isolate moiré mechanism \\\hline
Defect-Pinning Control & Declared grid & MoS2/WS2 with chalcogen vacancies & Sustained & Distinguish defect-mediated pinning from moiré potential barriers \\\hline
Contact-Resistance Control & Aligned (0$^\circ$) & MoS2/WS2 & Sustained & Detect thermal contact resistance drift mimicking intrinsic conductivity change \\\hline
Strain-Relaxation Control & Declared grid & MoS2/WS2 & Sustained & Identify wrinkle formation or localized delamination as alternative transport modifiers \\\hline
\end{tabular}
\end{table*}~\cite{cite_26ecaf73d88dfe33,cite_54b9175b8bb4c4e0,cite_78dc2ce387cb7c90,cite_a411cbbebf7838b9,cite_b9cc1fcb55e94349}

\subsection{Counterexample and Boundary Analysis}
The design boundary is explicit: structural stability of the engineered twist angle holds only when the interlayer binding energy exceeds the critical threshold set by the twist-angle-dependent moiré potential landscape and when defect-mediated or chemically-induced interlayer pinning is negligible. The declared falsifier is the relaxation of engineered twist angles to equilibrium stacking within device-relevant cycling conditions in the absence of defect or chemical pinning, accompanied by recovery of cross-plane conductivity to within 20\% of pristine aligned values. This falsifier satisfies all stated assumptions: it occurs within the declared cycling regime, excludes defect and chemical pinning by design of the control groups, and produces the predicted thermal signature of rotational relaxation. The design therefore distinguishes moiré-potential pinning from the four alternative explanations---interlayer diffusion, chemical bonding, contact-resistance drift, and strain relaxation---by requiring that the primary group's retention exceed the control groups' retention before any mechanism attribution is made.~\cite{cite_26ecaf73d88dfe33}

\subsection{Artifacts, Reproducibility, and Required Human Decisions}
Measurement instruments are specified by existing evidence: time-domain thermoreflectance for cross-plane thermal conductivity, Raman spectroscopy for layer identification and structural characterization, and scanning atomic electron tomography for atomic-coordinate precision at interfaces. Calibration and quality-control procedures remain unresolved and require human confirmation before execution. Reproducibility depends on the replication structure and batch definition for synthesis, transfer, and thermal cycling, which are also pending human input. The analysis plan, including randomization, missing-data handling, batch-effect treatment, repetition structure, and statistical model, is not finalized and must be confirmed by the researcher. Sampling eligibility, minimum measurable areas, and sample-size or power basis likewise await human specification. These items are procedural dependencies, not design failures; the protocol structure is complete, but its numerical parameters and statistical decision rules are gated on human review.~\cite{cite_26ecaf73d88dfe33,cite_9776eb140ec8227e,cite_b9cc1fcb55e94349}

\subsection{Transition to Next Stage}
The protocol, comparison groups, and boundary conditions are now specified. The next stage converts these design commitments into a preregistered decision matrix that maps prespecified outcome branches to allowed conclusions and next actions, ensuring that the analysis plan's statistical thresholds are fixed before any measurement is taken.

\section{Expected Outcome Branches and Conditional Conclusions}
\subsection{Baseline Mechanism and Operational Gap}
Twist-angle engineering in van der Waals heterostructures produces extreme thermal anisotropy by suppressing cross-plane phonon transport while preserving in-plane crystallinity. In random-stacked MoS2 and WS2 films, interlayer rotation yields a room-temperature thermal anisotropy ratio of approximately 900, with through-plane conductivity dropping to roughly 57 mW m$^{-1}$ K$^{-1}$ (W3202213959). This mechanism is validated only for those two materials, leaving applicability to other van der Waals families unverified. The persistence of engineered rotational configurations under sustained thermal cycling and operational stress gradients is likewise unaddressed. The present route replaces the assumed shear-strain locking mechanism with a twist-angle-dependent moiré potential landscape that creates energy barriers pinning specific twist angles. The central hypothesis H1 states that engineered angles remain structurally stable under thermal cycling provided interlayer binding energy exceeds a critical threshold determined by the local moiré potential, and in the absence of defect-mediated or chemically-induced interlayer pinning.~\cite{cite_26ecaf73d88dfe33}

\subsection{Decision Criteria and Competing Explanations}
Structural stability is accepted when the post-cycling twist angle remains within the pre-cycling uncertainty interval and cross-plane thermal conductivity retention exceeds the pre-specified threshold. The design must distinguish moiré-potential pinning from four declared alternatives: thermal cycling inducing interlayer diffusion or chemical bonding that pins layers independently of the moiré potential; changes in thermal contact resistance during cycling mimicking changes in intrinsic cross-plane conductivity; defect-mediated pinning via chalcogen vacancies maintaining suppressed conductivity; and strain relaxation via wrinkle formation or localized delamination altering thermal transport without actual twist-angle relaxation. A supporting branch requires the prespecified analysis to be consistent with the declared relation while planned controls do not favor any stated alternative. The falsifier is demonstration that engineered twist angles relax to equilibrium stacking within device-relevant cycling conditions in the absence of defect or chemical pinning, recovering cross-plane conductivity to within 20\% of pristine aligned values.~\cite{cite_26ecaf73d88dfe33,cite_6b49ace365c37cc5}

\begin{table*}[htbp]
\centering
\footnotesize
\setlength{\tabcolsep}{3pt}
\renewcommand{\arraystretch}{1.12}
\caption{Conditional Outcome Decision Matrix}
\label{tab:expected-outcomes-block-3}
\begin{tabular}{|p{0.23\textwidth}|p{0.23\textwidth}|p{0.23\textwidth}|p{0.23\textwidth}|}
\hline
Outcome Branch & Trigger Condition & Allowed Conclusion & Next Action \\\hline
supports\_mechanism & Prespecified analysis consistent with declared relation; controls do not favor stated alternative & Supports, but does not prove, moiré-potential pinning within design boundary & Replicate under independently confirmed conditions; test most consequential boundary condition \\\hline
partial\_or\_heterogeneous & Variation across declared conditions, units, or measurement contexts & Relation is conditional or heterogeneous; no universal conclusion warranted & Predefine and check plausible moderators; improve coverage of conditions and measurement comparability \\\hline
null\_or\_contradictory & Prespecified comparison does not support declared relation or favors alternative & Proposed relation not supported in this design boundary; absence of support is not proof of absence generally & Audit construct validity and comparison adequacy; revise mechanism or boundary claim before next iteration \\\hline
uninformative\_or\_invalid & Quality-control, missingness, protocol-deviation, or validity criteria prevent interpretation & No scientific conclusion warranted; planned design did not yield interpretable evidence & Resolve validity or data-quality failure before repeating design; obtain human confirmation of measurement, sampling, and analysis prerequisites \\\hline
\end{tabular}
\end{table*}

\subsection{Interpretive Boundaries and Material Generality}
The supports\_mechanism branch establishes conditional support only within the declared design boundary; it does not prove the moiré-potential threshold relation generally. The partial\_or\_heterogeneous branch signals that the relation depends on material family, cycling protocol, or measurement context, and no universal conclusion follows. The null\_or\_contradictory branch indicates that the proposed mechanism is not supported in this boundary, which is not equivalent to proof that moiré potential pinning never operates. The uninformative\_or\_invalid branch is a no-information condition: it means the planned design did not yield interpretable evidence, not that the mechanism failed or that the research plan is invalid. Material generality remains the most consequential open boundary. The extreme anisotropy mechanism is validated only for MoS2 and WS2, and whether random interlayer rotations yield similar anisotropy ratios in materials with different phonon dispersion or interlayer bonding characteristics is unverified. Mechanical evidence from boron nitride nanosheets shows that strong interlayer interactions prevent sliding under indentation (W2695145865), and van der Waals forces between stacked crystals give rise to layer-number dependent friction, superlubricity, and spatial modulation of elastic properties through moiré structures (W2803426300). These upstream findings constrain the scope of the moiré-potential mechanism without establishing it across families.~\cite{cite_6b49ace365c37cc5,cite_78dc2ce387cb7c90,cite_54b9175b8bb4c4e0}

\subsection{Transition to Risk and Limitations}
Each branch maps to a specific interpretive consequence and next action, preserving the separation between expected outcomes and observed results. No branch has been observed; all are prespecified decision rules for the proposed factorial design-of-experiments. The uninformative\_or\_invalid branch requires resolution of measurement, sampling, and analysis prerequisites before any scientific conclusion is drawn. The partial\_or\_heterogeneous branch demands predefined moderators and improved coverage across material families and cycling conditions. These interpretive boundaries feed directly into the risk and limitations stage, where human-review decisions on calibration standards, sample eligibility, and statistical thresholds determine whether the design can yield interpretable evidence in the first place.

\section{Risks, Limitations, and Human Review Requirements}
\subsection{Scope Boundaries and Mechanism Conflation}
The proposed relation is bounded to twist-angle-engineered thermal anisotropy in van der Waals heterostructures under operational thermal cycling, explicitly excluding regimes where defect-mediated or chemically-induced interlayer pinning dominates. This scope delimiter is necessary because the mechanism replacement hypothesis posits that the twist-angle-dependent moiré potential landscape creates fundamental energy barriers that pin specific twist angles, distinct from and dominant over defect-mediated or chemical pinning mechanisms. However, the upstream evidence for extreme thermal anisotropy via interlayer rotation is validated only for MoS2 and WS2; applicability to other van der Waals materials with different phonon dispersion or interlayer bonding characteristics remains unverified. Consequently, any generalization of the critical binding-energy threshold across material families is a limitation of the current design boundary, not an established property. The structural stability claim holds only where thermal cycling conditions minimize interlayer diffusion and chemical bonding, ensuring moiré potential pinning remains the dominant retention mechanism.~\cite{cite_26ecaf73d88dfe33,cite_6b49ace365c37cc5,cite_b9cc1fcb55e94349}

\begin{table*}[htbp]
\centering
\footnotesize
\setlength{\tabcolsep}{3pt}
\renewcommand{\arraystretch}{1.12}
\caption{Limitation and Human-Review Decision Matrix}
\label{tab:risk-limitations-and-review-rlr-02}
\begin{tabular}{|p{0.23\textwidth}|p{0.23\textwidth}|p{0.23\textwidth}|p{0.23\textwidth}|}
\hline
Limitation Domain & Evidence Status & Human-Review Decision & Release Condition Before Escalation \\\hline
Material Generality & Unverified beyond MoS2/WS2 & Confirm target van der Waals material families and justify cross-family extrapolation of moiré potential thresholds. & Signed material-scope addendum defining the eligible families and the basis for binding-energy comparability. \\\hline
Critical Threshold Definition & Equation or boundary condition for 'critical threshold' not provided & Specify the operational equation or measurable proxy for the twist-angle-dependent moiré potential energy barrier. & Approved threshold definition integrated into the hypothesis mapping. \\\hline
Thermal Cycling Parameters & Frequency, amplitude, and duration undefined & Define the sustained thermal-cycling protocol that suppresses diffusion and chemical bonding while imposing operational stress gradients. & Locked cycling protocol parameters with documented suppression rationale. \\\hline
Measurement Calibration & Instrument calibration standards and quality-control procedures undefined & Select and validate calibration standards for cross-plane thermal conductivity and twist-angle retention measurements. & Calibration plan with traceable standards and quality-control acceptance criteria. \\\hline
Statistical Decision Rule & Model and decision thresholds for stability support unspecified & Prespecify the statistical model, uncertainty interval, and conductivity-retention threshold governing the supports/null branches. & Finalized analysis plan with explicit decision thresholds. \\\hline
Replication and Batching & Replicate strategy and batch definitions for synthesis, transfer, and cycling undefined & Define batch structure and replication plan to enable moderator checks across material families and cycling intervals. & Approved replication and batch schema. \\\hline
Sample Eligibility & Eligibility criteria and minimum measurable areas unspecified & Establish minimum flake area and structural-integrity criteria compatible with the chosen characterization methods. & Eligibility criteria aligned with measurement sensitivity limits. \\\hline
\end{tabular}
\end{table*}~\cite{cite_26ecaf73d88dfe33}

\subsection{Formal and Evidentiary Scope Limits}
\paragraph{Human-review checklist.} The following scope limits are enforced by the authoring constraints and must be respected by all downstream sections: - Formal and empirical claims remain separate: the moiré potential energy-barrier mechanism is a proposed relation, not a proved theorem; no counterexample analysis applies outside a formal claim, and the current status is not\_run. - No observed results exist: all outcome branches are expected\_not\_observed; the supports\_mechanism branch would support but not prove the declared relation within the design boundary. - The uninformative\_or\_invalid branch is a procedural dependency: it indicates that prespecified quality-control, missingness, protocol-deviation, or validity criteria prevented interpretation, and it does not update the status of the mechanism claim nor render the research plan invalid. - Alternative explanations constrain interpretation: thermal contact resistance changes, defect-mediated pinning, thermally induced chemical bonding, and strain relaxation via wrinkle formation or localized delamination each mimic or mask the moiré potential effect; controls must be designed to exclude them before any supports\_mechanism conclusion is admissible. - The generalizability gap to materials beyond MoS2/WS2 is an open item, not a resolved limitation; it remains needs\_human\_input until the material-scope addendum is released.~\cite{cite_26ecaf73d88dfe33,cite_6b49ace365c37cc5,cite_b9cc1fcb55e94349}

\subsection{Concrete Human-Review Decisions and Release Conditions}
Before the design advances to execution, human review must resolve seven canonical needs\_human\_input records. The reviewer must (1) confirm the target van der Waals material families, because the extreme anisotropy mechanism is evidenced only for MoS2 and WS2 and the binding-energy threshold is material-dependent; (2) supply the numerical critical-threshold equation or measurable proxy, since the hypothesis's central conditional premise is currently undefined; (3) fix the thermal-cycling protocol parameters that suppress interlayer diffusion and chemical bonding, ensuring the moiré potential pinning mechanism is not confounded; (4) approve instrument calibration standards and quality-control procedures for both structural and thermal measurements, given that TDTR and twist-angle retention require distinct calibration regimes; (5) prespecify the statistical model and decision thresholds that map measured twist-angle retention and cross-plane conductivity retention to the supports\_mechanism, partial\_or\_heterogeneous, and null\_or\_contradictory branches; (6) define the replication structure and batch schema for synthesis, transfer, and cycling to enable the material-dependent binding-energy ablation; and (7) set sample eligibility criteria and minimum measurable areas compatible with the selected characterization methods. Each release condition above is the gate that converts a needs\_human\_input record into a design\_assumption or evidence\_backed premise. Until all seven are signed, the design remains in a pre-release state, and no outcome branch may be interpreted as supporting or refuting the central hypothesis.~\cite{cite_26ecaf73d88dfe33}

\section{Material System, Characterization, and DOE Plan}
\subsection{Material System and Structural Baseline}
The proposed material system consists of van der Waals heterostructures engineered with specific interlayer twist angles to suppress cross-plane thermal conductivity. The foundational premise relies on the established mechanism that interlayer rotational disorder impedes through-plane thermal transport while maintaining in-plane crystallinity, yielding a room-temperature thermal anisotropy ratio of approximately 900 in transition metal dichalcogenide films. This structural configuration is operationally defined as a stack of two or more 2D layers transferred to a designated substrate, with the target twist angle set during assembly. The experimental unit is a single heterostructure stack with a defined interlayer twist angle, and the design type is a factorial design-of-experiments. The selection of this material system is grounded in the evidence that random interlayer rotations break through-plane translational symmetry, creating a phonon-glass-like cross-plane transport regime. This baseline structural stability is the primary independent variable to be tested under sustained thermal cycling.~\cite{cite_26ecaf73d88dfe33}

\subsection{Measurement and Calibration Strategy}
\begin{itemize}
\item The characterization plan employs a multi-modal approach to isolate intrinsic thermal properties from interfacial artifacts and structural relaxation. The following protocols are designated for the pre-cycling baseline and post-cycling assessment:
\item **Thermal Conductivity Measurement:** Time-domain thermoreflectance (TDTR) is selected for measuring through-plane thermal conductivity. This technique requires a laser system with a 785 nm wavelength, 74.86 MHz repetition rate, and 18 mW total power, maintaining a steady-state temperature rise of less than 4 K. This specific calibration ensures that the measured signal reflects intrinsic cross-plane conductivity rather than transient heating artifacts.
\item **Structural Integrity and Twist Angle Retention:** Raman spectroscopy is utilized to identify and characterize the chemical and physical properties of the heterostructures. Specifically, layer breathing modes observed via resonant overtones in the 80--300 cm$^{-1}$ range will serve as a proxy for interlayer coupling strength and structural integrity. Twist-angle retention is operationally defined as the difference between the nominal assembly angle and the angle measured on the same region of the sample after a defined cycling interval.
\item **Control for Contact Resistance:** To address the alternative explanation that changes in thermal contact resistance mimic changes in intrinsic conductivity, measurement protocols will include controls for substrate coupling. The contact resistance between 2D materials and metals can be minimized to the 10$^{-7}$ $\Omega$ cm$^{2}$ range using specific metallization and post-annealing techniques, ensuring that the thermal bottleneck remains the interlayer interface rather than the electrode contact.
\end{itemize}~\cite{cite_26ecaf73d88dfe33,cite_c8ce7e25ee0b4157}

\begin{table*}[htbp]
\centering
\footnotesize
\setlength{\tabcolsep}{3pt}
\renewcommand{\arraystretch}{1.12}
\caption{Factorial Design of Experiments Matrix}
\label{tab:materials-and-characterization-doe-matrix}
\begin{tabular}{|p{0.23\textwidth}|p{0.23\textwidth}|p{0.23\textwidth}|p{0.23\textwidth}|}
\hline
Factor & Levels & Operational Definition & Primary Confounder Addressed \\\hline
**Interlayer Twist Angle** & Discrete Grid (e.g., 0$^\circ$, 5$^\circ$, 10$^\circ$, 15$^\circ$) & Nominal angle set during assembly. & Moiré potential depth variation \\\hline
**Material Family** & MoS$_{2}$/WS$_{2}$, h-BN, Black Phosphorus & Specific 2D layer composition. & Intrinsic binding energy differences \\\hline
**Thermal Cycling Intervals** & 0, 10, 50, 100 cycles & Repeated sequence of controlled temperature excursions. & Cumulative thermal stress \\\hline
**Defect Density** & Low, Medium, High & Chalcogen vacancy concentration. & Defect-mediated pinning \\\hline
**Substrate Type** & SiO$_{2}$, Sapphire, h-BN & Surface treatment and adhesion layer. & Interfacial thermal contact resistance \\\hline
\end{tabular}
\end{table*}~\cite{cite_26ecaf73d88dfe33,cite_b9cc1fcb55e94349}

\subsection{Scope Limitations and Human-Review Dependencies}
The proposed design assumes that interlayer binding energy and moiré potential depths can be accurately calculated or measured for the selected material families. However, the specific numerical targets for the twist-angle grid and the thermal-cycling protocol parameters (frequency, amplitude, duration) are currently undefined and require human input to establish device-relevant operational envelopes. Furthermore, the generalizability of the extreme thermal anisotropy mechanism beyond MoS$_{2}$ and WS$_{2}$ remains unverified, posing a risk to the material family factor in the DOE. The analysis plan, including statistical models for determining twist-angle stability and decision thresholds for accepting structural integrity, is marked as needing human input. Consequently, the transition to the next stage of research requires the resolution of these operational definitions to ensure that the factorial design can distinguish between moiré-potential-driven stability and defect-mediated or chemically-induced pinning.~\cite{cite_26ecaf73d88dfe33}

\appendices
\section{Idea Source Checkpoints and Direction Selection Audit}
The selected direction replaces the assumed shear-strain locking mechanism with a twist-angle-dependent moiré potential landscape as the primary determinant of thermal cycling stability in van der Waals heterostructures. This mechanism-replacement hypothesis posits that engineered interlayer twist angles remain structurally stable under sustained thermal cycling and operational stress gradients, provided the interlayer binding energy exceeds a critical threshold dictated by the local moiré potential energy barriers. The proposal explicitly bounds this stability to regimes free from defect-mediated or chemically-induced interlayer pinning, acknowledging that such confounding mechanisms could mask the true moiré potential threshold.~\cite{cite_26ecaf73d88dfe33}

\begin{table*}[htbp]
\centering
\footnotesize
\setlength{\tabcolsep}{3pt}
\renewcommand{\arraystretch}{1.12}
\caption{Decision matrix (Candidate)}
\label{tab:appendix-idea-evolution-checkpoint-comparison}
\begin{tabular}{|p{0.23\textwidth}|p{0.23\textwidth}|p{0.23\textwidth}|p{0.23\textwidth}|}
\hline
Checkpoint & Source File & Retained Elements & Scope Corrections \\\hline
Checkpoint 1 & idea\_candidate.json & Central hypothesis, claim scope, mechanism replacement & Establishes the baseline moiré potential landscape as the dominant retention mechanism distinct from defect pinning. \\\hline
Checkpoint 2 & idea\_result.json & Assumptions, boundary conditions, discriminating observation & Adds explicit assumptions regarding the accurate calculation of binding energies and the minimization of interlayer diffusion during cycling. \\\hline
\end{tabular}
\end{table*}~\cite{cite_26ecaf73d88dfe33}

The direction selection is anchored in the gap between established thermal anisotropy mechanisms and unverified operational stability. Prior research demonstrates that random interlayer rotations in transition metal dichalcogenide films impede through-plane thermal transport while maintaining in-plane crystallinity, yielding a room-temperature thermal anisotropy ratio of approximately 900 [cite\_26ecaf73d88dfe33]. However, this mechanism is validated exclusively for MoS2 and WS2, leaving its applicability to other van der Waals materials unverified. Furthermore, the persistence of these rotational configurations under sustained thermal cycling and operational stress gradients remains unaddressed, creating a critical missing assumption regarding structural stability under operation. The proposal integrates these boundary conditions by requiring material-dependent binding energy ablation and low-temperature cycling regimes to isolate the moiré potential mechanism from thermally induced chemical bonding.~\cite{cite_26ecaf73d88dfe33}

This audit preserves the open questions necessary for the subsequent experimental design phase. The exact equation or boundary condition for the critical threshold remains undefined, and the numerical targets for the interlayer twist-angle grid and thermal-cycling protocol require human input. Generalizability to materials beyond the initially validated MoS2 and WS2 systems is listed as a gap with an undefined scope. Consequently, the transition to the experimental design stage must prioritize the specification of synthesis and measurement protocols for controlled interlayer twist angles, as well as the determination of statistical models and decision thresholds required to evaluate the twist-angle stability hypothesis.~\cite{cite_26ecaf73d88dfe33}

\section{Variables, Symbols, and Operational Definitions}
\subsection{Operational Definitions and Measurement Anchors}
The experimental unit is a single van der Waals heterostructure stack with a defined interlayer twist angle. The primary dependent variable, cross-plane thermal conductivity, is anchored to time-domain thermoreflectance (TDTR), which measures through-plane conductivity using a 785 nm laser system with a 74.86 MHz repetition rate and 18 mW total power, maintaining a steady-state temperature rise below 4 K. Twist-angle retention is operationally defined as the difference between the nominal assembly angle and the angle measured on the same sample region after a defined cycling interval. Moiré structures spatially modulate elastic properties and interlayer friction, establishing the physical basis for the twist-angle-dependent potential landscape. Structural integrity is assessed via Raman spectroscopy, which identifies and characterizes the chemical and physical properties of the 2D materials.~\cite{cite_26ecaf73d88dfe33,cite_6b49ace365c37cc5,cite_9776eb140ec8227e}

\subsection{Confounding Variables and Boundary Conditions}
The central hypothesis posits that structural stability fails when the interlayer binding energy falls below a critical threshold determined by local moiré potential energy barriers, or when defect-mediated pinning and interlayer chemical bonding become the dominant retention mechanisms. Defect-mediated interlayer pinning and chemically-induced interlayer bonding are treated as confounders that mask the true moiré potential threshold. Interfacial thermal contact resistance is a critical control variable; changes during thermal cycling can mimic changes in intrinsic cross-plane thermal conductivity, necessitating rigorous separation of bulk and interface contributions. Strain relaxation via wrinkle formation or localized delamination alters thermal transport without actual twist-angle relaxation, requiring the structural integrity score to monitor for morphological degradation.

\begin{table*}[htbp]
\centering
\footnotesize
\setlength{\tabcolsep}{3pt}
\renewcommand{\arraystretch}{1.12}
\caption{Unresolved Operationalization Decisions}
\label{tab:appendix-variables-and-definitions-b3}
\begin{tabular}{|p{0.3067\textwidth}|p{0.3067\textwidth}|p{0.3067\textwidth}|}
\hline
Variable & Operational Definition Status & Required Human Input \\\hline
Controlled interlayer twist angles (V1) & Undefined synthesis and measurement protocol & Specify assembly technique and metrology for target angle setting \\\hline
Intrinsic cross-plane thermal conductivity (V2) & Measurement technique unspecified & Confirm TDTR or alternative method and calibration standards \\\hline
Critical threshold (V4) & Equation or boundary condition not provided & Define the mathematical threshold for binding energy failure \\\hline
Sustained thermal cycling (V5) & Parameters undefined & Specify frequency, amplitude, and duration of temperature excursions \\\hline
Material generalizability (V8) & Scope undefined & Determine applicability beyond MoS2/WS2 to other van der Waals families \\\hline
\end{tabular}
\end{table*}

\subsection{Transition to Experimental Design}
The operational definitions and confounding variables established in this appendix constrain the factorial design-of-experiments. The unresolved parameters for twist-angle grids, thermal-cycling protocols, and material families must be resolved to execute the pre-specified decision rule: structural stability is accepted if the post-cycling twist angle remains within the pre-cycling uncertainty interval and cross-plane thermal conductivity retention exceeds the pre-specified threshold. These dependencies directly inform the sampling and eligibility criteria for the subsequent experimental design.

\section{Evidence Coverage, Unknown Items, and Review Checklist}
\subsection{Evidence Coverage and Bounded Use of Sources}
The proposal's central premise---that twist-angle-dependent moiré potential landscapes dictate thermal cycling stability---rests on two pillars of upstream evidence: the demonstrated extreme anisotropy via interlayer rotation and the structural complexity of van der Waals interfaces. Kim et al. (W3202213959) established that random interlayer rotations in TMD films impede through-plane thermal transport while maintaining in-plane crystallinity, yielding a room-temperature thermal anisotropy ratio of approximately 900. This mechanism provides the baseline for the proposed structural stability hypothesis. However, the generalizability of this mechanism is strictly bounded; the extreme anisotropy was validated only for MoS2 and WS2, leaving applicability to other van der Waals materials with different phonon dispersion or interlayer bonding characteristics unverified. This gap necessitates the proposed factorial design-of-experiments to compare material families, but the specific selection of these families remains a design assumption pending human input. Furthermore, the interplay between moiré potentials and mechanical stability is supported by evidence that van der Waals forces give rise to phenomena such as superlubricity and spatial modulation of elastic properties through Moiré structures (W2803426300). Yet, a more thorough understanding of friction, fracture, and stress transfer mechanisms between 2D layers and with the substrate is still lacking, highlighting the need for the proposed operational stress gradients. The measurement of these properties relies on established techniques; time domain thermoreflectance (TDTR) was used to measure through-plane thermal conductivity in the foundational study (W3202213959), and Raman spectroscopy layer breathing modes can be observed in few-layer graphene via resonant overtones, scaling in frequency with layer number (W2024227324). These measurement capabilities are real, but their specific calibration standards for the proposed thermal cycling protocol are undefined.~\cite{cite_26ecaf73d88dfe33,cite_6b49ace365c37cc5}

\subsection{Review Checklist for Unresolved Design Parameters}
\begin{itemize}
\item The following items are critical dependencies for the experiment design. Each requires specific human input to resolve the ambiguity in the proposed methodology:
\item **Material System Selection (unknown-22)**: The specific van der Waals material families to be compared must be defined to address the gap in material generality. The current scope is undefined beyond MoS2/WS2.
\item **Design of Experiments (unknown-21)**: The factorial design structure must be specified, including the exact grid of interlayer twist angles and the number of thermal-cycling intervals.
\item **Process Variables (unknown-23)**: The parameters for 'sustained thermal cycling' (frequency, amplitude, duration) and the operational stress gradients must be numerically defined.
\item **Sample or Material Definition (unknown-25)**: Sample eligibility criteria and minimum measurable areas acceptable for the chosen characterization methods (e.g., TDTR, Raman) need to be established.
\item **Batches (unknown-20)**: The definition of synthesis, transfer, and thermal-cycling batches must be clarified to control for batch effects.
\item **Replicate Strategy (unknown-24)**: The planned replication structure and statistical decision thresholds for determining twist-angle stability must be specified.
\end{itemize}

\subsection{Release Criteria for Future Claims}
Future claims regarding the moiré potential threshold can only be released if the design resolves the confounding mechanisms identified in the reasoning context. Specifically, the structural stability hypothesis is bounded by the condition that defect-mediated or chemically-induced interlayer pinning is absent. If the proposed design does not explicitly control for these factors through the defined material systems and process variables, any observed retention of thermal anisotropy cannot be uniquely attributed to the moiré potential landscape. The decision rule for structural stability requires that the twist angle remains within the pre-cycling uncertainty interval and cross-plane thermal conductivity retention exceeds a pre-specified threshold. Until the critical threshold for interlayer binding energy is operationally defined and the specific material families are chosen, the mechanism replacement remains a proposed contribution rather than a testable design. The release of the 'supports\_mechanism' outcome branch is contingent upon the successful execution of the controls for interlayer diffusion and chemical bonding, which are currently listed as needs\_human\_input.\section*{References}
\bibliographystyle{IEEEtran}
\bibliography{references}
\end{document}
